\documentclass[conference]{IEEEtran}

\usepackage[utf8]{inputenc}
\usepackage[T1]{fontenc}
\usepackage[brazil]{babel}
\usepackage{graphicx}
\usepackage{booktabs}
\usepackage{multirow}
\usepackage{subcaption}
\usepackage{amsmath, amssymb}
\usepackage{float}
\usepackage{xcolor}
\usepackage{fancybox}

\graphicspath{{assets/titanic/plots/}{assets/titanic/tables/}}

% Placeholder para tabelas (substituir por screenshot)
\newcommand{\tablePlaceholder}[1]{%
  \begin{center}
    \fbox{\parbox{0.85\columnwidth}{%
      \centering\vspace{0.4cm}
      \textbf{[PLACEHOLDER]} \\[0.2cm]
      \textit{#1} \\[0.4cm]
    }}
  \end{center}
}

% -------------------------------------------------------------------------
\title{Análise Estatística Descritiva do Dataset Titanic:\\
Perfil Socioeconômico e Padrões de Sobrevivência}

\author{
  \IEEEauthorblockN{Heitor Teixeira}
  \IEEEauthorblockA{
    \textit{MBA: Ciência de dados}\\
    UNIFOR\\
    Fortaleza, Brasil\\
    heitor.mba.unifor@gmail.com
  }
}

\begin{document}

\maketitle

% -------------------------------------------------------------------------
\begin{abstract}
Este trabalho apresenta uma análise estatística descritiva do dataset
Titanic, com foco na caracterização socioeconômica dos passageiros e nos
padrões de sobrevivência observados no naufrágio de 1912. O dataset
utilizado compreende 872 registros de passageiros após filtragem de tarifas
nulas. As tarifas originais em libras esterlinas de 1912 foram corrigidas
monetariamente para valores de 2024 por meio de um índice composto, que
combina dados do Banco da Inglaterra (1912--2016) e do Banco Mundial
(2016--2024), resultando em fator acumulado de 107,14. A análise inclui
classificação das variáveis, estatísticas descritivas completas, percentis,
distribuições de frequência, análise de densidade por classe, perfis de
sobrevivência por gênero, classe e faixa etária, além de análise de
correlação de Pearson entre idade e tarifa. Os resultados evidenciam que a
sobrevivência foi fortemente determinada por fatores sociais e demográficos:
passageiros de primeira classe apresentaram taxa de sobrevivência de 63\%,
contra 24\% na terceira classe; mulheres sobreviveram a uma taxa de 73\%,
enquanto homens apresentaram apenas 19\%. A correlação entre idade e tarifa
mostrou-se estatisticamente significativa porém praticamente negligenciável
($r = 0{,}115$; $p = 0{,}0007$), indicando independência entre essas
variáveis.
\end{abstract}

\begin{IEEEkeywords}
Estatística descritiva, Titanic, análise exploratória de dados,
correlação de Pearson, visualização de dados.
\end{IEEEkeywords}

% =========================================================================
\section{Introdução}

Este trabalho realiza uma análise de estatística descritiva sobre o dataset
Titanic (D1), em atendimento à atividade avaliativa da disciplina de
Estatística Descritiva. O dataset contém registros de passageiros com
variáveis demográficas, socioeconômicas e de sobrevivência, constituindo
um conjunto adequado para a aplicação das técnicas exigidas.

A análise foi implementada em Python com as bibliotecas \textit{pandas},
\textit{matplotlib}, \textit{numpy}, \textit{scipy} e \textit{requests},
e cobre os seguintes requisitos: (i) classificação das variáveis qualitativas
e quantitativas; (ii) gráficos de barras e pizza para Sex, Pclass e
Survived; (iii) medidas de tendência central e dispersão para Age e Fare;
(iv) quartis e percentis P10 e P90 de Age; (v) histograma de Age; e
(vi) gráfico de dispersão e coeficiente de Pearson entre Age e Fare.

Como a atividade não especificou um formato de entrega, optou-se pelo
formato de artigo científico segundo o padrão IEEE, por ser uma estrutura
consolidada na comunicação técnica e científica, que favorece a organização
clara dos resultados e sua reprodutibilidade.

% =========================================================================
\section{Fundamentação Teórica}

\subsection{Estatística Descritiva}

A estatística descritiva constitui o conjunto de técnicas que visa sumarizar
e apresentar de forma sistemática as características de um conjunto de dados.
Seu objetivo mas sim organizar e descrever os dados de forma clara e objetiva.

\subsection{Medidas de Posição e Dispersão}

As medidas de posição central mais utilizadas são a média aritmética,
a mediana e a moda. A média aritmética é sensível a valores extremos
(\textit{outliers}), enquanto a mediana oferece maior robustez em
distribuições assimétricas. As medidas de dispersão, como o desvio-padrão e
o intervalo interquartil (IQR), quantificam a variabilidade dos dados ao
redor de uma medida central.

Os percentis dividem o conjunto de dados ordenado em cem partes iguais.
Os quartis $Q_1$ (P25), $Q_2$ (P50) e $Q_3$ (P75) são os percentis de maior
uso prático. Neste trabalho, calculam-se também os percentis P10 e P90,
que delimitam as caudas da distribuição e complementam a caracterização
da variabilidade da variável Age.

\subsection{Distribuições de Frequência e Histogramas}

Um histograma representa a distribuição de frequências de uma variável
contínua por meio de retângulos adjacentes, onde a base corresponde a um
intervalo de classe e a altura à frequência (ou densidade) observada. O
número de classes pode ser determinado pela regra da raiz quadrada
($k \approx \sqrt{n}$), adotada neste trabalho, conforme falado em sala de aula.

A estimativa de densidade (\textit{Kernel Density Estimation},
KDE) é uma abordagem não paramétrica para estimar a função de densidade de
probabilidade de uma variável aleatória. Neste trabalho,
a KDE gaussiana foi utilizada para sobrepor curvas de densidade contínua
sobre os histogramas e para comparar distribuições de tarifas entre classes.

\subsection{Correlação de Pearson}

O coeficiente de correlação de Pearson $r$ mede a intensidade e a direção
da relação linear entre duas variáveis contínuas, assumindo valores em
$[-1, 1]$:
\begin{equation}
  r = \frac{\sum_{i=1}^{n}(x_i - \bar{x})(y_i - \bar{y})}
           {\sqrt{\sum_{i=1}^{n}(x_i - \bar{x})^2 \cdot
                  \sum_{i=1}^{n}(y_i - \bar{y})^2}}
\end{equation}
Valores de $|r|$ próximos de 1 indicam forte relação linear; valores próximos
de 0 indicam ausência de relação linear. O teste de significância associado
avalia a hipótese nula $H_0: r = 0$.

\subsection{Correção Monetária}

A comparação de valores monetários ao longo do tempo requer a aplicação de
índices de inflação. Neste trabalho, as tarifas originais em libras de 1912
foram corrigidas para 2024 utilizando dados históricos do Banco da Inglaterra
e do Banco Mundial, possibilitando uma interpretação contemporânea do poder
aquisitivo dos passageiros.

% =========================================================================
\section{Metodologia}

\subsection{Dataset e Filtragem}

O dataset utilizado é o clássico conjunto de dados do Titanic, contendo
informações dos passageiros a bordo no momento do naufrágio. Por decisão
do autor, registros com tarifa igual a zero foram removidos: sem valor de
tarifa registrado, não é possível confirmar que o indivíduo era um passageiro
pagante, e sua inclusão distorceria as medidas de tendência central e
dispersão da variável Fare. A amostra final utilizada nas análises compreende
\textbf{872 passageiros}. As variáveis foram renomeadas para o português e
mapeadas conforme descrito na Tabela~\ref{tab:variaveis}.

\begin{figure}[H]
  \centering
  \includegraphics[width=\columnwidth]{table_3_2.png}
  \caption{Identificação e classificação das variáveis do dataset Titanic.}
  \label{tab:variaveis}
\end{figure}

\subsection{Correção Monetária das Tarifas}

As tarifas originais estão expressas em libras esterlinas de 1912. Para
permitir uma interpretação contemporânea, foi construído um índice de
correção composto a partir de duas fontes:

\begin{itemize}
  \item \textbf{Banco da Inglaterra (BoE)}: fator acumulado de 1912 a 2016
        igual a $81{,}46\times$.
  \item \textbf{Banco Mundial (WB)}: fator acumulado de 2016 a 2024 igual a
        $1{,}3153\times$.
\end{itemize}

O fator total resultante é $107{,}14 = 81{,}46 \times 1{,}3153$, aplicado
multiplicativamente a cada tarifa original para obter a coluna
\textit{Tarifa 2024}. A composição detalhada do índice é apresentada na
Tabela~\ref{tab:inflacao}.

\begin{figure}[H]
  \centering
  \includegraphics[width=0.9\columnwidth]{table_2_2_1.png}
  \caption{Composição do índice de correção monetária (1912→2024).}
  \label{tab:inflacao}
\end{figure}

\subsection{Faixas Etárias}

A variável idade contínua foi discretizada em quatro faixas etárias
definidas pelo autor com base em critérios demográficos intuitivos,
conforme a Tabela~\ref{tab:faixas}:

\begin{figure}[H]
  \centering
  \includegraphics[width=0.65\columnwidth]{table_2_4.png}
  \caption{Distribuição de passageiros por faixa etária.}
  \label{tab:faixas}
\end{figure}

As quatro faixas definidas foram: \textbf{Criança} (0--12 anos, $n=79$),
\textbf{Jovem} (12--17 anos, $n=51$), \textbf{Adulto} (17--59 anos,
$n=711$) e \textbf{Idoso} (59--100 anos, $n=31$).

% =========================================================================
\section{Resultados}

\subsection{Estatísticas Descritivas}

A Tabela~\ref{tab:descritivas} apresenta as principais estatísticas
descritivas das variáveis quantitativas do dataset.

\begin{figure}[H]
  \centering
  \includegraphics[width=\columnwidth]{table_3_1.png}
  \caption{Estatísticas descritivas das variáveis quantitativas do dataset Titanic.}
  \label{tab:descritivas}
\end{figure}

Os principais resultados são:

\begin{itemize}
  \item \textbf{Idade}: média de 29,43 anos ($\sigma = 14{,}19$), mediana
        de 28 anos; amplitude total de 0,42 a 80 anos.
  \item \textbf{Tarifa 2024}: média de £3.520,61 ($\sigma = £5.359,63$),
        mediana de £1.553,47; máximo de £54.888,84. A razão entre média e
        mediana ($\approx 2{,}27$) evidencia forte assimetria positiva.
\end{itemize}

\subsection{Distribuição de Frequências: Tarifa}

A Figura~\ref{fig:hist_fare} exibe o histograma das tarifas corrigidas para
2024, com eixo $x$ em escala logarítmica para acomodar a amplitude da
distribuição. O número de classes foi definido pela regra da raiz quadrada
($k \approx \sqrt{872} \approx 29$).

\begin{figure}[H]
  \centering
  \includegraphics[width=\columnwidth]{hist_fare.png}
  \caption{Distribuição de frequência da tarifa por passageiro (valores
           corrigidos para 2024, escala logarítmica). Titanic.}
  \label{fig:hist_fare}
\end{figure}

A distribuição é fortemente assimétrica à direita: a grande maioria dos
passageiros concentra-se em tarifas inferiores a £2.000, enquanto um pequeno
grupo pagou valores superiores a £20.000.

\subsection{Distribuição de Frequências: Idade}

A Figura~\ref{fig:hist_age} apresenta o histograma da idade com curva de
densidade por kernel sobreposta.

\begin{figure}[H]
  \centering
  \includegraphics[width=\columnwidth]{hist_age.png}
  \caption{Distribuição de frequência da idade por passageiro com
           estimativa de densidade por kernel (KDE gaussiana). Titanic.}
  \label{fig:hist_age}
\end{figure}

Ao contrário da tarifa, a distribuição de idades apresenta comportamento
aproximadamente simétrico (média $\approx$ mediana $\approx$ 28--29 anos),
com concentração entre 20 e 40 anos e cauda à direita menos pronunciada.
A coorte infantil (abaixo de 12 anos) é visível mas representa parcela
reduzida da amostra.

\subsection{Distribuição das Variáveis Qualitativas}

As Figuras~\ref{fig:gender}, \ref{fig:class} e~\ref{fig:survival} exibem,
respectivamente, a distribuição de gênero, classe e sobrevivência, cada uma
representada por um gráfico de barras e um gráfico de pizza.

\begin{figure}[H]
  \centering
  \includegraphics[width=\columnwidth]{gender.png}
  \caption{Distribuição de gênero: contagem e proporção. Titanic.}
  \label{fig:gender}
\end{figure}

\begin{figure}[H]
  \centering
  \includegraphics[width=\columnwidth]{class.png}
  \caption{Distribuição por classe: contagem e proporção. Titanic.}
  \label{fig:class}
\end{figure}

\begin{figure}[H]
  \centering
  \includegraphics[width=\columnwidth]{survival.png}
  \caption{Distribuição de sobrevivência: contagem e proporção. Titanic.}
  \label{fig:survival}
\end{figure}

O conjunto de passageiros é composto predominantemente por homens ($\approx
65\%$, $n=568$) e passageiros de terceira classe ($\approx 55\%$, $n=480$).
A taxa geral de sobrevivência é de 38\% ($n=335$), o que indica que a
maioria dos passageiros ($\approx 62\%$, $n=537$) não sobreviveu ao
naufrágio.

\subsection{Densidade da Tarifa por Classe}

A Figura~\ref{fig:density} apresenta as curvas KDE das tarifas
(escala logarítmica) estratificadas por classe, evidenciando o grau de
diferenciação socioeconômica entre os grupos.

\begin{figure}[H]
  \centering
  \includegraphics[width=\columnwidth]{density.png}
  \caption{Densidade da tarifa por classe (KDE, escala logarítmica). Titanic.}
  \label{fig:density}
\end{figure}

A terceira classe apresenta distribuição unimodal concentrada ao redor de
£900; a segunda classe exibe comportamento bimodal com picos em £1.500 e
£3.000, sugerindo subdivisões internas de acomodação; a primeira classe
possui cauda direita longa, atingindo valores próximos a £55.000, e
variabilidade interna significativamente maior.

\subsection{Perfil de Sobrevivência}

\subsubsection{Sobrevivência por Classe}

\begin{figure}[H]
  \centering
  \includegraphics[width=\columnwidth]{survival_class.png}
  \caption{Sobrevivência por classe (barras empilhadas proporcionais). Titanic.}
  \label{fig:survival_class}
\end{figure}

Como ilustrado na Figura~\ref{fig:survival_class}, a sobrevivência foi
fortemente estratificada pela classe socioeconômica:

\begin{itemize}
  \item \textbf{1ª Classe}: 63\% de sobrevivência.
  \item \textbf{2ª Classe}: 47\% de sobrevivência.
  \item \textbf{3ª Classe}: 24\% de sobrevivência.
\end{itemize}

A diferença de 39 pontos percentuais entre a primeira e a terceira classe
é a maior variação observada entre grupos no dataset.

\subsubsection{Sobrevivência por Gênero}

\begin{figure}[H]
  \centering
  \includegraphics[width=\columnwidth]{survival_gender.png}
  \caption{Sobrevivência por gênero (barras empilhadas proporcionais). Titanic.}
  \label{fig:survival_gender}
\end{figure}

A Figura~\ref{fig:survival_gender} evidencia o maior efeito observado na
análise: mulheres apresentaram taxa de sobrevivência de 73\%, enquanto
homens apresentaram apenas 19\%, uma diferença de 54 pontos percentuais.

\subsubsection{Sobrevivência por Faixa Etária}

\begin{figure}[H]
  \centering
  \includegraphics[width=\columnwidth]{survival_group_age.png}
  \caption{Sobrevivência por faixa etária (barras empilhadas proporcionais). Titanic.}
  \label{fig:survival_group_age}
\end{figure}

A Figura~\ref{fig:survival_group_age} demonstra que crianças obtiveram as
maiores taxas de sobrevivência entre as faixas etárias. Adultos, apesar de
constituírem a maioria da amostra ($n=711$), apresentaram as menores taxas
relativas.

\subsection{Boxplot e Percentis}

A Figura~\ref{fig:boxplot} exibe os boxplots de idade e tarifa 2024, com
eixo $y$ em escala logarítmica para a tarifa a fim de acomodar os valores
extremos.

\begin{figure}[H]
  \centering
  \includegraphics[width=\columnwidth]{boxplot.png}
  \caption{Boxplot: Idade (esquerda) e Tarifa 2024 em escala logarítmica
           (direita). Titanic.}
  \label{fig:boxplot}
\end{figure}

Os percentis calculados para cada variável são apresentados nas
Tabelas~\ref{tab:percentis_age} e~\ref{tab:percentis_fare}.

\begin{figure}[H]
  \centering
  \includegraphics[width=0.55\columnwidth]{table_3_8_2.png}
  \caption{Percentis da variável Idade (P10, Q1, Q2, Q3, P90).}
  \label{tab:percentis_age}
\end{figure}

\begin{figure}[H]
  \centering
  \includegraphics[width=0.6\columnwidth]{table_3_8_3.png}
  \caption{Percentis da variável Tarifa 2024 (P10, Q1, Q2, Q3, P90).}
  \label{tab:percentis_fare}
\end{figure}

Para a idade, os quartis Q1, Q2 e Q3 correspondem a 20, 28 e 38 anos,
respectivamente. Para a tarifa 2024, o Q1 está em £849,05, a mediana em
£1.553,47 e o Q3 em £3.350,67, com P90 atingindo £8.352,13, evidenciando
a forte concentração de tarifas nas faixas mais baixas.

\subsection{Correlação entre Idade e Tarifa}

A Figura~\ref{fig:scatter} apresenta o diagrama de dispersão entre idade e
tarifa 2024, com os pontos coloridos por sobrevivência, e a
Figura~\ref{fig:correlacao} exibe a matriz de correlação de Pearson.

\begin{figure}[H]
  \centering
  \includegraphics[width=\columnwidth]{scatter_pearson.png}
  \caption{Dispersão: Idade $\times$ Tarifa 2024. Coeficiente de Pearson
           $r = 0{,}115$. Titanic.}
  \label{fig:scatter}
\end{figure}

\begin{figure}[H]
  \centering
  \includegraphics[width=\columnwidth]{matriz_correlacao.png}
  \caption{Matriz de correlação de Pearson (Idade $\times$ Tarifa 2024).
           Titanic.}
  \label{fig:correlacao}
\end{figure}

O coeficiente de Pearson foi calculado via função \texttt{pearsonr} da
biblioteca \textit{scipy.stats}, que retorna o valor de $r$ e o
\textit{p-value} associado. O valor obtido foi $r = 0{,}115$ ($p = 0{,}0007$). Embora
estatisticamente significativo ao nível de 0,1\%, o valor indica uma
correlação positiva \emph{muito fraca} sob qualquer escala de referência
convencional. Os resultados da correlação são detalhados na
Tabela~\ref{tab:pearson}.

\begin{figure}[H]
  \centering
  \includegraphics[width=0.6\columnwidth]{table_3_9.png}
  \caption{Coeficiente de correlação de Pearson entre Idade e Tarifa 2024.}
  \label{tab:pearson}
\end{figure}

Esse resultado implica que a idade dos passageiros não determinou o preço
da passagem: as tarifas foram definidas pela classe de acomodação escolhida,
e não por características demográficas individuais. Do ponto de vista de
modelagem preditiva, essa independência é favorável, pois indica ausência
de multicolinearidade entre essas duas variáveis.

% =========================================================================
\section{Conclusão}

Este trabalho conduziu uma análise estatística descritiva completa do dataset
Titanic, investigando o perfil dos passageiros e os determinantes da
sobrevivência no naufrágio de 1912.

A análise revelou que a composição da amostra é majoritariamente masculina
($\approx 65\%$) e de terceira classe ($\approx 55\%$), refletindo o perfil
imigrante típico das travessias transatlânticas da época. A distribuição de
tarifas é fortemente assimétrica à direita, com mediana de £1.553 e média
de £3.521 (valores 2024), enquanto a distribuição de idades é aproximadamente
simétrica, centrada nos 28--29 anos.

Os resultados mais expressivos foram encontrados na análise de sobrevivência:
a taxa variou de 24\% (terceira classe) a 63\% (primeira classe), e de
19\% (homens) a 73\% (mulheres). Esses dados sugerem que a sobrevivência
não foi aleatória: a diferença de 39 pontos percentuais entre classes pode
indicar acesso preferencial aos botes e proximidade dos camarotes ao convés;
a diferença de 54 pontos entre gêneros é consistente com um protocolo de
evacuação que priorizou mulheres e crianças, grupo que também apresentou
as maiores taxas por faixa etária.

Por fim, a correlação de Pearson entre idade e tarifa ($r = 0{,}115$)
demonstrou que, embora estatisticamente significativa, a relação linear entre
essas variáveis é praticamente inexistente, indicando que o preço da
passagem foi determinado pela classe de acomodação e não pela idade do
passageiro.


\end{document}
